\section{Применение к данным}
\label{sec:results}

\subsection{Прямая задача и опубликованное решение}
\draft{Для Модели~1 сравнение выполнено на исходной геометрии XLSX по значениям на кровле, подошве и средним по слоям. На рис.~\ref{fig:data1-temperature} и~\ref{fig:data1-ro} показаны современные поля; знаковая относительная разность равна $100(u^{\mathrm{calc}}-u^{\mathrm{ref}})/u^{\mathrm{ref}}$. В картах ошибок T используется температура в градусах Цельсия. Для сводных оценок используется табличная норма без пространственных весов:}

Для Модели~1 предоставлено эталонное решение прямой задачи, с которым сопоставляются результаты разработанной программы. Сравнение выполняется на исходной геометрии по значениям температуры и отражательной способности на кровле, подошве и в среднем по мощности каждого слоя. Современные поля приведены на рис.~\ref{fig:data1-temperature} и~\ref{fig:data1-ro}. Знаковая относительная разность определяется как $100(u^{\mathrm{calc}}-u^{\mathrm{ref}})/u^{\mathrm{ref}}$, поэтому положительные значения соответствуют превышению расчёта над эталоном. При нормировании температурных различий используются градусы Цельсия.

Для общей оценки различий используется относительная евклидова норма по сопоставляемым табличным значениям. Индекс $j$ перебирает записи выбранного поля и типа значения, например средние по слоям за все конфигурации. Каждая запись учитывается с одинаковым весом, без умножения на площадь соответствующего участка разреза. Поэтому эта оценка характеризует согласование выходных таблиц и отличается от нормы по площади, использованной при сравнении полей в синтетической задаче

\begin{equation}
 \varepsilon_{\mathrm{tab}}(u) = 100
 \frac{\left(\sum_j (u_j^{\mathrm{calc}} - u_j^{\mathrm{ref}})^2\right)^{1/2}}
 {\left(\sum_j (u_j^{\mathrm{ref}})^2\right)^{1/2}}\quad[\%].
\end{equation}
\reportfigure{14_data1_temperature.png} {Сравнение современных температурных полей Модели~1. (а) предоставленное эталонное решение; (б) решение разработанной программы; (в) знаковая относительная разность расчёта и эталона. Температура выражена в градусах Цельсия. На панелях (а) и (б) используется общая цветовая шкала.} {fig:data1-temperature}
\reportfigure{15_data1_reflectance.png} {Сравнение современных полей отражательной способности витринита Модели~1. (а) предоставленное эталонное решение; (б) решение разработанной программы; (в) знаковая относительная разность расчёта и эталона. На панелях (а) и (б) используется общая цветовая шкала значений $R_o$.} {fig:data1-ro}
\draft{По средним значениям слоёв ошибки T/Ro за все 18~возрастов равны 2,803/1,859\%; для современности~--- 3,598/2,385\%. На профилях (рис.~\ref{fig:profiles}) полоса показывает $\pm2\%$ от эталона как масштаб сравнения, а не заданную погрешность эталонного решения. Для Модели~2 опубликованного эталона нет; пример ускорения за счёт упрощения геометрии приведён на рис.~\ref{fig:rdp} и в табл.~\ref{tab:rdp}.}

По средним значениям слоёв за все 18~возрастов относительная норма различий с эталоном составляет 2,803\% для температуры и 1,859\% для отражательной способности. Для современного состояния получены 3,598 и 2,385\% соответственно. Эти числа оценивают близость результатов двух программ. Численная точность предложенного метода определяется отдельно по задачам с известным решением, рассмотренным в разделе~\ref{sec:verification}.

Распределение различий по глубине показано на рис.~\ref{fig:profiles} для сечений $x=282$~м и $x=1400$~м. Серая полоса соответствует отклонению на $\pm2\%$ от эталонного значения и используется как общий масштаб визуального сравнения. Она не является оценкой известной погрешности эталонного решения. Для Модели~2 эталон не предоставлен, поэтому этот набор использован для исследования влияния упрощения геометрии на время и результат расчёта (рис.~\ref{fig:rdp}, табл.~\ref{tab:rdp}).

\reportfigure{16_data1_profiles.png} {Современные профили температуры и отражательной способности в Модели~1. (а), (б) температура при $x=282$~м и $x=1400$~м; (в), (г) отражательная способность в тех же сечениях. Чёрные линии обозначают эталонное решение, красные~--- расчёт. Серая полоса соответствует отклонению на $\pm2\%$ от эталона и служит масштабом сравнения.} {fig:profiles}
\FloatBarrier

\subsection{Обратная задача}
\draft{Полные поля при найденном потоке показаны на рис.~\ref{fig:inverse-fields}. Границы слоёв и скважины позволяют сопоставить расчёт с расположением наблюдений.}

В обратной задаче Модели~3 базальный поток определяется по 500~измерениям температуры и 50~измерениям отражательной способности витринита в скважинах~A--E. На рис.~\ref{fig:inverse-fields} показаны современные поля, полученные при восстановленном потоке, а также границы слоёв и положения скважин. Такое представление позволяет соотнести температурный режим и накопленную зрелость с геологическим строением разреза. Степень согласования с измерениями далее оценивается по профилям и невязкам отдельно для каждой скважины.

\reportfigure{17_inverse_fields.png} {Современные поля Модели~3 при восстановленном базальном потоке. (а) температура; (б) отражательная способность витринита. На разрезах показаны границы геологических слоёв и положения скважин~A--E, по измерениям в которых решалась обратная задача.} {fig:inverse-fields}
\draft{Начальное приближение~--- первый фактически рассчитанный поток после сдвига параметров внутрь допустимой области. Он сопоставлен с найденным на рис.~\ref{fig:inverse-q}. Изменение потока на последней эпохе требует дальнейшего исследования весов и размерности параметризации.}

На рис.~\ref{fig:inverse-q} сопоставлены начальная и восстановленная истории базального потока. Начальная история соответствует первому набору параметров, для которого выполнен прямой расчёт, после сдвига начального приближения внутрь допустимого множества. Подбираются все 17~безразмерных параметров, а преобразование в узловые значения потока обеспечивает его монотонность. В восстановленной истории заметно изменение потока на последней эпохе. Для выяснения его устойчивости требуется исследовать влияние числа параметров и весов наблюдений. На основании одного найденного решения нельзя считать эту особенность установленной характеристикой истинного потока.

\reportfigure{18_inverse_flux.png} {История базального потока в обратной задаче Модели~3. Синяя линия обозначает начальное приближение, красная~--- восстановленный поток. Точками отмечены значения на возрастах конфигураций, между которыми выполняется линейная интерполяция.} {fig:inverse-q}
\draft{Подбор выполнен на сетке $200\times100$~м: 10~итераций, 44~прямых вычисления, $J=265{,}883$; остановка по лимиту итераций. На рис.~\ref{fig:optimization} верхняя шкала отмечает концы итераций; нулевая итерация включает 18 начальных вычислений. Контроль на сетке $200\times50$~м даёт $J_T=205{,}898$, $J_{R_o}=55{,}608$.}

При подборе использованы шаги $\Delta x=200$~м, $\Delta z=100$~м и $\Delta t=1$~млн~лет. За 10~итераций выполнены 44~прямых расчёта, а минимальное полученное значение функционала составило $J=265{,}883$. Вычисления завершены по заданному ограничению числа итераций. На рис.~\ref{fig:optimization} показаны значения функционала при каждом прямом расчёте и минимальное значение среди уже рассмотренных наборов параметров. Верхняя шкала отмечает завершение итераций оптимизационного метода. Начальная итерация включает 18~прямых расчётов, необходимых для построения приближений компонент невязки.

Для проверки результата найденный поток использован в дополнительном прямом расчёте с более мелким вертикальным шагом $\Delta z=50$~м при тех же $\Delta x=200$~м и временном шаге. Получены значения $J_T=205{,}898$ и $J_{R_o}=55{,}608$, в сумме $J=261{,}506$. Параметры потока на этой сетке повторно не подбирались. Все последующие профили и оценки невязок относятся к данному контрольному расчёту.

\reportfigure{19_inverse_objective.png} {Изменение целевой функции при решении обратной задачи Модели~3. Чёрные точки обозначают значения для пробных наборов параметров, красная линия~--- минимальное значение среди уже рассмотренных наборов. Нижняя шкала показывает число прямых расчётов, верхняя~--- завершённые итерации оптимизационного метода.} {fig:optimization}
\draft{Профили при начальном и найденном потоках показаны на рис.~\ref{fig:wells-temperature} и~\ref{fig:wells-ro}; полоса соответствует заданной погрешности измерений $\pm2\%$. Оба расчёта выполнены на сетке $200\times50$~м. В допуске 456/500~температур и 39/50~значений Ro. RMS относительных невязок~--- 1,283\% и 2,109\%, максимумы~--- 6,743\% и 6,907\% соответственно.}

На рис.~\ref{fig:wells-temperature} и~\ref{fig:wells-ro} показаны профили температуры и отражательной способности при начальном и восстановленном потоках вместе с измерениями и интервалами их погрешности $\pm2\%$. Оба прямых расчёта выполнены на сетке $200\times50$~м, чтобы исключить влияние различий сетки при сравнении профилей. После подбора в заданные интервалы попадают 456 из 500~температур и 39 из 50~значений отражательной способности. Среднеквадратичные относительные невязки равны 1,283 и 2,109\%, а максимальные относительные невязки~--- 6,743 и 6,907\% соответственно.

В табл.~\ref{tab:wells} приведены число согласованных наблюдений и среднеквадратичные относительные невязки для каждой скважины. Наибольшая температурная невязка по этой норме составляет 2,017\% в скважине~A. Для отражательной способности наибольшие значения получены в скважинах~D и E и равны 2,565 и 3,652\%. Таким образом, восстановленный поток улучшает согласование с наблюдениями, но не обеспечивает попадания всех расчётных значений в заданные интервалы. Среднеквадратичная относительная невязка $\varepsilon_{\mathrm{rms}}$ для группы из $N_{\mathrm{obs}}$~наблюдений определяется выражением
\begin{equation}
 \varepsilon_{\mathrm{rms}}(u) = 100\left[
 \frac{1}{N_{\mathrm{obs}}}\sum_j
 \left(\frac{u_j^{\mathrm{calc}} - u_j^{\mathrm{obs}}}
 {u_j^{\mathrm{obs}}}\right)^2\right]^{1/2}\quad[\%],
 \qquad u \in \{T_C,R_o\}.
\end{equation}

\reportfigure{20_wells_temperature.png} {Сравнение расчётных и измеренных температур в скважинах Модели~3. (а)--(д) скважины~A--E соответственно. Синие линии показывают расчёт при начальном потоке, красные~--- при восстановленном потоке, чёрные точки~--- измерения. Серая полоса соответствует принятой погрешности измерений $\pm2\%$.} {fig:wells-temperature}
\reportfigure{21_wells_reflectance.png} {Сравнение расчётной и измеренной отражательной способности витринита в скважинах Модели~3. (а)--(д) скважины~A--E соответственно. Синие линии показывают расчёт при начальном потоке, красные~--- при восстановленном потоке, чёрные точки~--- измерения. Серая полоса соответствует принятой погрешности измерений $\pm2\%$.} {fig:wells-ro}
\begin{table}[htbp]
\centering\small
\renewcommand{\arraystretch}{1.12}
\caption{Невязки и число согласованных наблюдений по скважинам.}\label{tab:wells}
\begin{tabular}{lrrrr}
\toprule
\multirow{2}{*}{Скважина} & \multicolumn{2}{c}{$T$} & \multicolumn{2}{c}{$R_o$} \\
\cmidrule(lr){2-3}\cmidrule(lr){4-5}
 & В пределах 2\% & $\varepsilon_{\mathrm{rms}}$, \% & В пределах 2\% & $\varepsilon_{\mathrm{rms}}$, \% \\
\midrule
A & 79/100 & 2,080 & 10/10 & 1,073 \\
B & 94/100 & 0,979 & 10/10 & 1,111 \\
C & 94/100 & 1,073 & 10/10 & 0,466 \\
D & 95/100 & 0,922 & 8/10 & 2,458 \\
E & 95/100 & 1,058 & 3/10 & 3,433 \\
\bottomrule
\end{tabular}
\end{table}

\draft{Чувствительность результата при фиксированном найденном потоке приведена в табл.~\ref{tab:sensitivity}. Базовые параметры: $\Delta x=200$~м, $\Delta z=50$~м, $\Delta t=1$~млн~лет, $\varepsilon=0{,}05$~м; изменяется один параметр за расчёт.}

Влияние численных параметров проверено при фиксированном восстановленном потоке. За основу принят расчёт с $\Delta x=200$~м, $\Delta z=50$~м, $\Delta t=1$~млн~лет и $\varepsilon=0{,}05$~м. В каждом дополнительном расчёте изменяется только один параметр, а различия оцениваются в точках наблюдений (табл.~\ref{tab:sensitivity}). Уменьшение вертикального шага до 25~м изменяет температуру не более чем на 0,162\%, отражательную способность~--- на 0,147\%. При уменьшении временного шага до 0,25~млн~лет соответствующие изменения не превышают 0,00160 и 0,1191\%. Изменение минимальной высоты в проверенном диапазоне оказывает ещё меньшее влияние.

Изменения решения при проверенных сгущениях существенно меньше наибольших оставшихся невязок измерений. Следовательно, в рассмотренных пределах уменьшение пространственного и временного шагов не устраняет расхождения в скважинах. Для их объяснения необходимо отдельно исследовать подбор параметров и чувствительность результата к принятым предположениям о тепловой истории среды.

\begin{table}[htbp]
\centering\small
\caption{Максимальные относительные изменения полей в точках наблюдений при изменении параметров дискретизации.}\label{tab:sensitivity}
\begin{tabular}{lrrr}
\toprule
Параметр & Значение & $T$, \% & $R_o$, \% \\
\midrule
$\Delta z$, м & 25 & 0,16112 & 0,14695 \\
$\Delta z$, м & 100 & 0,25356 & 0,39488 \\
$\Delta t$, млн~лет & 0,25 & 0,00160 & 0,11908 \\
$\Delta t$, млн~лет & 0,5 & 0,00133 & 0,11930 \\
$\varepsilon$, м & 0,01 & 0,00070 & 0,01585 \\
$\varepsilon$, м & 0,1 & 0,00087 & 0,02004 \\
\bottomrule
\end{tabular}
\end{table}

